\scanpage{23}
\textcircled{\small IV} Où on laisse tomber la condition de fidélité.

Considérons des données (11) satisfaisant (13) mais pas néc. (12) ; donc à part l'unimodularité, on laisse tomber la condition $\varphi$ fidèle. Que se passe-t-il ?

\uline{Proposition 4} \inserted{\text{Sous les conditions précédentes}} \deleted{Pour que $\varphi$ soit fidèle, il faut et}

(a) Si $v\in D_n^0$ i.e. $v=u^i$, alors $v_E=\mathrm{id}\iff v(s_0)=s_0$.

(b) Pour que $\varphi$ soit fidèle, il f et s. que $\varphi\mid D_n^0$ soit fidèle, \deleted{i.e. que} \inserted{\text{ou encore}} (par a)) que $u^i s_0=s_0\Longrightarrow i\equiv0\ (n)$ (i.e. $u^i=1$), ou encore que $\forall d\in{\mathbb N}^*$, $d\mid n$, on ait $u^d(s_0)\ne s_0$ \unsure{$N$|$\mathcal N$}

(c) \deleted{Soit} Le noyau de $\varphi$ \deleted{\unsure{$G$}} est \deleted{un sous-groupe de} engendré par un (unique) élément $u^{\deleted{\unsure{i}}d}$, où \deleted{$i$ est} $d$ est un diviseur de $n$ tel que \deleted{$d\geq3$. \unsure{--} de sorte que $\varphi$ \unsure{est induit une repr. fi}} \inserted{\text{$n'=n/d\geq3$}}

\uline{Dém} (a) $v_E=\mathrm{id}\iff v(s_{\nu})=s_\nu\ \forall\nu$ (car les $s_\nu$ engend\textsuperscript{t} aff\textsuperscript{t} $E$) or $s_\nu=u^\nu s_0$ donc $v(s_\nu)=vu^\nu s_0=u^\nu v s_0$, donc il suffit $v(s_0)=s_0$.
% archive sheet number: 22
